% GENERATED FILE. Edit canonical lessons, then run npm run generate:books.
% canonical-source-hash: fnv1a64:6adb3e76eaa421b9
% canonical-lessons: PA-C297-mulaim, PA-C297-barik, PA-C297-kula, PA-C297-khundha, PA-C297-chubhna

\chapter{Rough and Smooth}
\label{ch:pa-a2-rough-and-smooth}
\hlchaptermodality{297}

\begin{chapteropening}
\textbf{By the end of this chapter:} \emph{I can say how something feels in the hand.}

Complete the last lesson of chapter 297: I can say how something feels in the hand.
\end{chapteropening}

\section[\texorpdfstring{\pa{ਮੁਲਾਇਮ} (mulāim)}{ਮੁਲਾਇਮ (mulāim)}]{\pa{ਮੁਲਾਇਮ} (mulāim) — smooth; soft to the touch}

\label{lesson:PA-C297-mulaim}



\begin{quote}
Before the new one: say the Punjabi for slick, then the Punjabi for to melt.
\end{quote}

\subsection*{You'll want to know: \pa{ਮੁਲਾਇਮ}}
\textbf{\pa{ਮੁਲਾਇਮ}} — \emph{mulāim} — \textquotedblleft{}smooth; soft to the touch\textquotedblright{}.

\begin{quote}
\pa{ਇਹ} \pa{ਕੱਪੜਾ} \pa{ਬਹੁਤ} \pa{ਮੁਲਾਇਮ} \pa{ਹੈ।}
\end{quote}

\begin{quote}
\emph{eh kappṛā bahut mulāim hai.}
\end{quote}

\begin{quote}
This cloth is very smooth.
\end{quote}

\begin{cousinweb}
\textbf{\pa{ਮੁਲਾਇਮ}} (\emph{mulāim}) is a loan, usually traced through Persian to Arabic \emph{mulāyim}, \textquotedblleft{}gentle, soft\textquotedblright{}.
\end{cousinweb}

\begin{grammarlens}[title={What you've built}]
One more word for this chapter.
\end{grammarlens}

\subsection*{Your turn}
\begin{itemize}
\raggedright
  \item \emph{Say it:} \emph{mulāim}
  \item \emph{Say it:} \emph{mulāim}, once more
  \item \emph{Say it:} \emph{eh kappṛā bahut mulāim hai.} — This cloth is very smooth.
  \item \emph{Recall:} say the Punjabi for slick, then the Punjabi for to melt, then say \emph{mulāim} again
\end{itemize}

\begin{tcolorbox}[breakable,colback=teal!4,colframe=teal!35!black,title={Before you move on}]
What does \pa{ਮੁਲਾਇਮ} mean? (\textquotedblleft{}Smooth; soft to the touch\textquotedblright{}.) Say it once more.
\end{tcolorbox}
\section[\texorpdfstring{\pa{ਬਰੀਕ} (barīk)}{ਬਰੀਕ (barīk)}]{\pa{ਬਰੀਕ} (barīk) — fine; fine-grained}

\label{lesson:PA-C297-barik}



\begin{quote}
Before the new one: say the Punjabi for to melt, then the Punjabi for smooth.
\end{quote}

\subsection*{You'll want to know: \pa{ਬਰੀਕ}}
\textbf{\pa{ਬਰੀਕ}} — \emph{barīk} — \textquotedblleft{}fine; fine-grained\textquotedblright{}.

\begin{quote}
\pa{ਦਰਿਆ} \pa{ਦੇ} \pa{ਕਿਨਾਰੇ} \pa{ਬਰੀਕ} \pa{ਰੇਤ} \pa{ਹੈ।}
\end{quote}

\begin{quote}
\emph{dariā de kināre barīk ret hai.}
\end{quote}

\begin{quote}
There is fine sand on the bank of the river.
\end{quote}

\begin{cousinweb}
\textbf{\pa{ਬਰੀਕ}} (\emph{barīk}) is a loan, usually traced to Persian \emph{bārīk}, \textquotedblleft{}thin, fine\textquotedblright{}.
\end{cousinweb}

\begin{grammarlens}[title={What you've built}]
One more word for this chapter.
\end{grammarlens}

\subsection*{Your turn}
\begin{itemize}
\raggedright
  \item \emph{Say it:} \emph{barīk}
  \item \emph{Say it:} \emph{barīk}, once more
  \item \emph{Say it:} \emph{dariā de kināre barīk ret hai.} — There is fine sand on the bank of the river.
  \item \emph{Recall:} say the Punjabi for to melt, then the Punjabi for smooth, then say \emph{barīk} again
  \item \emph{Write it:} copy \textbf{\pa{ਬਰੀਕ}} once, checking the whole word
\end{itemize}

\begin{tcolorbox}[breakable,colback=teal!4,colframe=teal!35!black,title={Before you move on}]
What does \pa{ਬਰੀਕ} mean? (\textquotedblleft{}Fine; fine-grained\textquotedblright{}.) Say it once more.
\end{tcolorbox}
\section[\texorpdfstring{\pa{ਕੂਲਾ} (kūlā)}{ਕੂਲਾ (kūlā)}]{\pa{ਕੂਲਾ} (kūlā) — tender; soft}

\label{lesson:PA-C297-kula}



\begin{quote}
Before the new one: say the Punjabi for smooth, then the Punjabi for fine-grained.
\end{quote}

\subsection*{You'll want to know: \pa{ਕੂਲਾ}}
\textbf{\pa{ਕੂਲਾ}} — \emph{kūlā} — \textquotedblleft{}tender; soft\textquotedblright{}.

\begin{quote}
\pa{ਬੱਚੇ} \pa{ਦੇ} \pa{ਹੱਥ} \pa{ਕੂਲੇ} \pa{ਹੁੰਦੇ} \pa{ਹਨ।}
\end{quote}

\begin{quote}
\emph{bachche de hatth kūle hunde han.}
\end{quote}

\begin{quote}
A child's hands are soft.
\end{quote}

\begin{cousinweb}
\textbf{\pa{ਕੂਲਾ}} (\emph{kūlā}) is inherited from Sanskrit \emph{kōmala}, \textquotedblleft{}tender, soft\textquotedblright{}; CDIAL lists Punjabi \emph{kūlā} under it (CDIAL 3523).
\end{cousinweb}

\begin{grammarlens}[title={What you've built}]
One more word for this chapter.
\end{grammarlens}

\subsection*{Your turn}
\begin{itemize}
\raggedright
  \item \emph{Say it:} \emph{kūlā}
  \item \emph{Say it:} \emph{kūlā}, once more
  \item \emph{Say it:} \emph{bachche de hatth kūle hunde han.} — A child's hands are soft.
  \item \emph{Recall:} say the Punjabi for smooth, then the Punjabi for fine-grained, then say \emph{kūlā} again
  \item \emph{Write it:} copy \textbf{\pa{ਕੂਲਾ}} once, checking the whole word
\end{itemize}

\begin{tcolorbox}[breakable,colback=teal!4,colframe=teal!35!black,title={Before you move on}]
What does \pa{ਕੂਲਾ} mean? (\textquotedblleft{}Tender; soft\textquotedblright{}.) Say it once more.
\end{tcolorbox}
\section[\texorpdfstring{\pa{ਖੁੰਢਾ} (khuṇḍhā)}{ਖੁੰਢਾ (khuṇḍhā)}]{\pa{ਖੁੰਢਾ} (khuṇḍhā) — blunt}

\label{lesson:PA-C297-khundha}



\begin{quote}
Before the new one: say the Punjabi for fine-grained, then the Punjabi for tender.
\end{quote}

\subsection*{You'll want to know: \pa{ਖੁੰਢਾ}}
\textbf{\pa{ਖੁੰਢਾ}} — \emph{khuṇḍhā} — \textquotedblleft{}blunt\textquotedblright{}.

\begin{quote}
\pa{ਇਹ} \pa{ਚਾਕੂ} \pa{ਖੁੰਢਾ} \pa{ਹੈ।}
\end{quote}

\begin{quote}
\emph{eh chākū khuṇḍhā hai.}
\end{quote}

\begin{quote}
This knife is blunt.
\end{quote}

\begin{cousinweb}
\textbf{\pa{ਖੁੰਢਾ}} (\emph{khuṇḍhā}) is inherited from a reconstructed \emph{*khuṇṭha}, \textquotedblleft{}defective\textquotedblright{}; CDIAL lists Punjabi \emph{khuṇḍhā}, \textquotedblleft{}blunt\textquotedblright{}, under it (CDIAL 3899).
\end{cousinweb}

\begin{grammarlens}[title={What you've built}]
One more word for this chapter.
\end{grammarlens}

\subsection*{Your turn}
\begin{itemize}
\raggedright
  \item \emph{Say it:} \emph{khuṇḍhā}
  \item \emph{Say it:} \emph{khuṇḍhā}, once more
  \item \emph{Say it:} \emph{eh chākū khuṇḍhā hai.} — This knife is blunt.
  \item \emph{Recall:} say the Punjabi for fine-grained, then the Punjabi for tender, then say \emph{khuṇḍhā} again
  \item \emph{Write it:} copy \textbf{\pa{ਖੁੰਢਾ}} once, checking the whole word
\end{itemize}

\begin{tcolorbox}[breakable,colback=teal!4,colframe=teal!35!black,title={Before you move on}]
What does \pa{ਖੁੰਢਾ} mean? (\textquotedblleft{}Blunt\textquotedblright{}.) Say it once more.
\end{tcolorbox}
\section[\texorpdfstring{\pa{ਚੁਭਣਾ} (chubhṇā)}{ਚੁਭਣਾ (chubhṇā)}]{\pa{ਚੁਭਣਾ} (chubhṇā) — to prick; to stick into}

\label{lesson:PA-C297-chubhna}



\begin{quote}
Before the new one: say the Punjabi for tender, then the Punjabi for blunt.
\end{quote}

\subsection*{You'll want to know: \pa{ਚੁਭਣਾ}}
\textbf{\pa{ਚੁਭਣਾ}} — \emph{chubhṇā} — \textquotedblleft{}to prick; to stick into\textquotedblright{}.

\begin{quote}
\pa{ਸੂਈ} \pa{ਉਂਗਲ} \pa{ਵਿੱਚ} \pa{ਚੁਭਦੀ} \pa{ਹੈ।}
\end{quote}

\begin{quote}
\emph{sūī uṅgal vich chubhdī hai.}
\end{quote}

\begin{quote}
The needle pricks the finger.
\end{quote}

Past (it …): \pa{ਉਹ} \pa{ਸੂਈ} \pa{ਚੁਭੀ}, \emph{oh sūī chubhī}.

\begin{cousinweb}
\textbf{\pa{ਚੁਭਣਾ}} (\emph{chubhṇā}) is inherited from a reconstructed \emph{*cubhyatē}, \textquotedblleft{}is pierced\textquotedblright{}; CDIAL lists Punjabi \emph{cubhṇā}, \textquotedblleft{}to be pierced\textquotedblright{}, under it (CDIAL 4867).
\end{cousinweb}

\begin{grammarlens}[title={What you've built}]
One more word for this chapter.
\end{grammarlens}

\subsection*{Your turn}
\begin{itemize}
\raggedright
  \item \emph{Say it:} \emph{chubhṇā}
  \item \emph{Say it:} \emph{chubhṇā}, once more
  \item \emph{Say it:} \emph{sūī uṅgal vich chubhdī hai.} — The needle pricks the finger.
  \item \emph{Recall:} say the Punjabi for tender, then the Punjabi for blunt, then say \emph{chubhṇā} again
  \item \emph{Write it:} copy \textbf{\pa{ਚੁਭਣਾ}} once, checking the whole word
\end{itemize}

\begin{tcolorbox}[breakable,colback=teal!4,colframe=teal!35!black,title={Before you move on}]
What does \pa{ਚੁਭਣਾ} mean? (\textquotedblleft{}To prick; to stick into\textquotedblright{}.) Say it once more.
\end{tcolorbox}
